\chapter{Conclusiones y Trabajo Futuro}

% --------------------------------------------------------
\section{Conclusiones}
% --------------------------------------------------------
% TODO: ~400 palabras — responder directamente cada pregunta de investigación
%
% 1. Sí, el ML (XGBoost, Grupo C) supera al ARIMA: RMSE 49.01 vs 107.48 pb (reducción 54.4%)
%
% 2. El NLP aporta: 2.30-2.59 pb de mejora incremental. La señal es de volumen
%    de cobertura mediática, no de sentimiento.
%
% 3. Granger: el NLP no anticipa el EMBI con rezago. La relación es contemporánea;
%    el EMBI anticipa el tono de las noticias (no al revés).
%
% 4. Quiebres estructurales confirmados en 2020-03 y 2022-01. El EMBI ecuatoriano
%    es exógeno en períodos normales, parcialmente endógeno en crisis.
%
% 5. El autoencoder confirma que la representación NLP es esencialmente lineal —
%    resultado negativo informativo que guía trabajo futuro.

% --------------------------------------------------------
\section{Implicaciones para política económica}
% --------------------------------------------------------
% TODO: ~200 palabras
% - Un sistema de monitoreo del EMBI en tiempo real basado en este modelo
%   podría alertar al Ministerio de Finanzas sobre presiones de costo de deuda
% - La dominancia post-COVID del DXY y Treasury 10Y implica que la política
%   fiscal ecuatoriana tiene poco margen de acción sobre el EMBI en el corto plazo
% - El desempleo local como factor COVID sugiere que la credibilidad fiscal
%   interna sí importa en escenarios de estrés extremo

% --------------------------------------------------------
\section{Trabajo futuro}
% --------------------------------------------------------
% TODO: ~300 palabras — lista priorizada
%
% 1. LSTM end-to-end: red recurrente sobre secuencia completa (macro + NLP) — podría
%    capturar dependencias temporales de orden superior que XGBoost ignora
%
% 2. NLP con texto crudo: embeddings FinBERT sobre artículos de noticias GDELT
%    (requiere acceso a texto, no solo métricas agregadas). Hipótesis: el sentimiento
%    del texto crudo sí anticipa el EMBI cuando el tono agregado no lo hace
%
% 3. Extender a 2008 con GDELT v1: incluir el default Correa requiere reconciliar
%    dos versiones del dataset con dimensiones diferentes
%
% 4. Modelos por régimen: entrenar XGBoost separado para cada subperiodo (pre-COVID,
%    COVID, post-COVID) y comparar con el modelo global
%
% 5. Panel de países EM: extender el análisis a Ecuador + Colombia + Perú + Argentina
%    para identificar qué factores son idiosincráticos de Ecuador vs regionales
